\subsection*{Цель работы}
Ознакомление с основными приемами работы с документацией при
составлении программ для микроконтроллеров. Приобретение навыков работы с
осциллографом и оценочной платой MCBSTM32F200 в качестве измерительного
генератора.

\subsection*{Постановка задачи}
Используя библиотеки Keil µVision5, разработать
программу для микроконтроллера (МК) STM32F200, которая включает и выключает
светодиод.

\subsection*{Методика выполнения работы}

В среде Keil µVision5 был создан проект, где была разработана программа, включающая и выключающая
светодиод. Программа использует переменные из библиотеки, упрощающие работу с ркгистрами МК STM32F200. На листинге \ref{lis:prog-code} представлен код программы.

\begin{lstlisting}[language=C, caption={Прогамма упавления светодиодом.}, label={lis:prog-code}, numbers=left]
#include "stm32f2xx.h"

void delay () {
    unsigned long i; // Counter for blinky delay
    i=0;
    for(i=0; i<2000000; i++) {} // Delay
}

int main () {
    RCC->AHB1ENR |= 1ul<<6; // Enable port G clocking
    GPIOG->MODER = (GPIOG->MODER & ~1ul<<15) | 1ul<<14;
    for (;;) {
        GPIOG->ODR |= 1ul<<7;
        delay ();
        GPIOG->ODR &= ~1ul<<7;
        delay ();
    }
}
\end{lstlisting}

\begin{itemize}
    \item Функция \verb|delay()| устанавливает задержку с помощью использования цикла \verb|for(i=0; i<2000000; i++) {}|, чтобы было видно мигание светодиода.
    \item В функции \verb|main()| включается тактирование порта G и устанавливается его режим работы; устанавливается режим выхода для порта PG7.
    \item Бесконечный цикл \verb|for(;;)| обеспецивает включение и выключение светодиода PG7 с периодичностью, заданной в функции \verb|delay()|.
\end{itemize}

\subsection*{Работа с осциллографом}

На фотографиях \ref{fig:IMG_20260310_122836.jpg} и \ref{fig:IMG_20260310_122844.jpg} показано подключение щупов осциллографа к плате.

\screenshot{IMG_20260310_122844.jpg}{Подкоючение щупа осциллографа к выходу PG7.}
\screenshot{IMG_20260310_122836.jpg}{Подключение заземления осциллографа к плате.}

Далее на фотографии \ref{fig:IMG_20260310_122848_1.jpg} изображён сам осциллограф.
По данным, отображённым на экране прибора, были получены следующие результаты:

\screenshot{IMG_20260310_122848_1.jpg}{Осциллограф, показывающий сигнал ШИМ.}

\begin{itemize}
    \item \textbf{Период сигнала} $T = 500\,\text{мс} \cdot 3 = 1500\,\text{мс} = 1.5\,\text{с}$;
    \item \textbf{Размах сигнала} $V_{pp} \approx 3.2\,\text{В}$ (масштаб 2 В/дел, высота $\approx 1.6$ деления);
    \item \textbf{Частота сигнала} $f = \frac{1}{T} = \frac{1}{1.5} \approx 0.67\,\text{Гц}$;
    \item \textbf{Длительность положительного уровня} $T_{\text{вкл}} = 650\,\text{мс}$;
    \item \textbf{Длительность отрицательного уровня} $T_{\text{выкл}} = 850\,\text{мс}$;
    \item \textbf{Коэффициент заполнения} $D = \frac{T_{\text{вкл}}}{T} = \frac{650}{1500} \approx 43\,\%$;
    \item \textbf{Время нарастания фронта} $t_r \approx 5.36\,\text{нс}$ (измерено при развертке 5 нс/дел).
\end{itemize}

На фотографии \ref{fig:IMG_20260310_125133.jpg} показаны данные прибора для измерения времени нарастания.

\screenshot{IMG_20260310_125133.jpg}{Данные осциллографа для измерения времени нарастания.}
\subsection*{Результаты работы}

Таким образом, в ходе выполнения лабораторной работы был разработан проект в среде Keil $\mu$Vision5 и написана программа на языке C для микроконтроллера STM32F200.

В процессе работы были освоены следующие навыки:
\begin{itemize}
    \item Работа с документацией и доступом к регистрам микроконтроллера.
    \item Работа с измерительным оборудованием: подключение щупов осциллографа и снятие характеристик сигнала.
\end{itemize}

На выводе PG7 был сгенерирован сигнал. С помощью цифрового осциллографа были получены его основные параметры: период $T=1.5$\,с, частота $f \approx 0.67$\,Гц и коэффициент заполнения $D \approx 43\,\%$. Также было измерено время нарастания фронта, которое составило $t_r \approx 5.36$\,нс, что подтверждает корректную работу выходного каскада микроконтроллера.
